\documentclass[10pt,a4paper]{amsart}
\usepackage[margin=19mm]{geometry}
\usepackage{amsmath,amssymb,mathtools}
\usepackage[scr=rsfs,cal=euler]{mathalfa}
\usepackage{xfrac}
\usepackage[unicode,hidelinks,bookmarks=false]{hyperref}
\newcommand{\ie}{i.e.\ }
\newcommand{\cf}{cf.\ }
\newcommand{\resp}{respectively\ }
\newcommand{\Subsection}{\subsection}
\providecommand{\nobreakdash}{\nobreak\hbox{-}}
\setlength{\emergencystretch}{2em}
\multlinegap0pt
\usepackage{amssymb}
\newtheorem{lemma}{Lemma}
\newtheorem{proposition}{Proposition}
\title{Truncations of the ring of number-theoretic functions, revisited}
\author{Jan Snellman}
\date{Source: arXiv:2608.29555v1 (CC BY 4.0)}
\begin{document}
\maketitle

\noindent\emph{Source and licence.} Truncations of the ring of number-theoretic functions, revisited. Version arXiv:2608.29555v1 (CC BY 4.0), distributed by arXiv under the Creative Commons Attribution 4.0 International licence, http://creativecommons.org/licenses/by/4.0/, as recorded in the arXiv licence metadata for this exact version and shown on the abstract page https://arxiv.org/abs/2608.29555v1. Full licence notice: \texttt{licenses/arxiv-2608.29555-CC-BY-4.0.txt}.\par\smallskip

\noindent\textit{Editorial scope.} Counted unit: the article's proposition prop:T (source0.tex lines 1020--1026). The article's complete original proof of it is delivered with it (source0.tex lines 1028--1069, one \texttt{proof} environment of 42 lines). No mathematics is written, repaired or re-derived: the statement and the proof are the article's own lines, in the article's own notation, and no retained line is substituted or re-flowed.  Retained uncounted as the source-local context the proof needs: lines 942--949; lines 984--995.  Nothing was omitted from any retained span, so the package carries no omission record.  No retained line contains a citation, so the wrapper carries no bibliography and no bibliography entry.  No retained line contains a figure environment, an image-inclusion command or drawing code, so no image is delivered and the package declares no asset dependency.  Editorial formatting only, in addition: an AMS \texttt{amsart} preamble that restates the article's own declarations and macros used by the retained lines, namely the environments \texttt{lemma}, \texttt{proposition} on separate counters, together with the standard packages those lines need.  The retained environments print in this document as Lemma 1, Proposition 1 in order of appearance, whereas the article prints the same items with the numbering of its own full document.  The complete native source of the article as distributed in the arXiv e-print for this exact version is bundled unchanged as \texttt{sources/208865/source0.tex}.
  \begin{equation}\label{eq:pnt-sub}
    \begin{aligned}
      \pi\bigl(Xe^{-a}\bigr)
        &= \frac{X}{M}\left(e^{-a}+\theta_1\,\frac{A_1(W+1)}{M}\right),\\
      \pi\bigl(Xe^{a}\bigr)
        &= \frac{X}{M}\left(e^{a}+\theta_2\,\frac{A_1(W+1)e^{W}}{M}\right),
    \end{aligned}
  \end{equation}

\begin{lemma}\label{lem:tail}
  There is an absolute constant \(c_0\) such that, for every \(W\) with
  \(1 \le W \le M/2\),
  \[
    \sum_{p \le Xe^{-W}}\pi(p)\,\pi(n/p)
      \ \le \ c_0\,e^{-W}\,\frac{X^{3}}{M^{3}}
      \ +\ c_0\,\frac{X^{3}}{M^{3}} \cdot \frac{\log n}{n^{1/4}} .
  \]
  The same bound holds for \(\sum_{p \le Xe^{-W}}g(p)\) and for
  \(\sum_{p \le Xe^{-W}}(\pi(p)-1)\pi(n/p)\), both summands being at most
  \(\pi(p)\pi(n/p)\).
\end{lemma}

\begin{proposition}\label{prop:T}
  There are absolute constants \(C_0\) and \(n_0\) such that for \(n \ge n_0\)
  \[
    \left|\;\frac{M^{3}}{X^{3}}\,T(n)-1\;\right|
      \ \le \ C_0\,\frac{\log M}{\sqrt M} .
  \]
\end{proposition}

\begin{proof}
  Fix \(W \ge 1\) with \(M \ge 2W\), to be chosen. Split \(T\) at \(Xe^{-W}\); by
  Lemma~\ref{lem:tail} the part with \(p \le Xe^{-W}\) is at most
  \(c_0(e^{-W}+\log n/n^{1/4})X^{3}/M^{3}\).

  In the main range \(Xe^{-W} < p \le X\) put \(a = \log(X/p) \in [0,W)\). By
  \eqref{eq:pnt-sub},
  \begin{align*}
    \bigl(\pi(p)-1\bigr)\pi(n/p)
      =\frac{X^{2}}{M^{2}}
       &\Bigl(e^{-a}+\theta_1\tfrac{A_1(W+1)}{M}-\tfrac{M}{X}\Bigr)\\
       &\times
       \Bigl(e^{a}+\theta_2\tfrac{A_1(W+1)e^{W}}{M}\Bigr).
  \end{align*}
  Expanding, the leading term is \(e^{-a}e^{a} = 1\), and every other term
  carries a factor \(A_1(W+1)e^{W}/M\) or \(M/X\); using \(e^{-a} \le 1\),
  \(e^{a} \le e^{W}\) and \(X \ge M^{2}\), all of them together are at most
  \(A_2(W+1)e^{W}/M\) in absolute value, with \(A_2\) absolute. Hence
  \begin{align*}
    \sum_{Xe^{-W} < p \le X}\bigl(\pi(p)-1\bigr)\pi(n/p)
      ={}&\frac{X^{2}}{M^{2}}
          \Bigl(1+\theta_3\,\tfrac{A_2(W+1)e^{W}}{M}\Bigr)\\
        &\times \#\{\,p : Xe^{-W} < p \le X\,\},
  \end{align*}
  \(|\theta_3|\le 1\) --- this is where the summand being constant to leading
  order does the work: no partition of \([0,W]\) and no Riemann sum is needed,
  only the number of primes in the range. That number is, by
  \eqref{eq:pnt-sub} again,
  \[
    \pi(X)-\pi\bigl(Xe^{-W}\bigr)
      =\frac{X}{M}\Bigl(1-e^{-W}+\theta_4\,\tfrac{2A_1(W+1)}{M}\Bigr).
  \]
  Multiplying, and adding the tail bound,
  \[
    \frac{M^{3}}{X^{3}}T(n)
      =1+O\!\left(e^{-W}+\frac{(W+1)e^{W}}{M}+\frac{\log n}{n^{1/4}}\right),
  \]
  with an absolute implied constant. Now take \(W = \tfrac{1}{2}\log M\), legitimate
  for \(M \ge e^{2}\) since then \(W \ge 1\) and \(M \ge 2W\). It gives
  \(e^{-W} = M^{-1/2}\) and \((W+1)e^{W}/M = (\tfrac{1}{2}\log M+1)/\sqrt M\), both
  \(O(\log M/\sqrt M)\), while \(\log n/n^{1/4}\) is smaller than either.
\end{proof}

\end{document}
